\documentclass[aspectratio=169,10pt]{beamer}
\usetheme{Madrid}
\usecolortheme{default}
\usepackage[utf8]{inputenc}
\usepackage[T1]{fontenc}
\usepackage{lmodern}
\usepackage{amsmath,amssymb}
\usepackage{enumitem}
\usepackage{xcolor}

\setbeamertemplate{navigation symbols}{}
\setbeamertemplate{footline}[frame number]
\setbeamertemplate{itemize items}[circle]
\setbeamerfont{frametitle}{size=\large,series=\bfseries}

\title{CIA Part 1 -- Questions 355--413}
\subtitle{CAE--Board Relations, Independence Threats, QAIP, Conformance \& External Assessment}
\author{CIA Study Prep}
\date{\today}

\begin{document}
	
	\begin{frame}
		\titlepage
	\end{frame}
	
	\begin{frame}{Instructions}
		\begin{itemize}
			\item 59 questions in two domains only: (A) CAE--board relations and independence threats; (B) QAIP, conformance judgment, and external assessment.
			\item Each question slide is followed by a separate answer/explanation slide.
			\item Every option is defensible; the correct answer turns on a precise Standards priority or a single word.
			\item Read every qualifier (MOST, LEAST, FIRST, BEST, PRIMARY) before answering.
		\end{itemize}
	\end{frame}
	
	% ============================================================
	% DOMAIN A: CAE--BOARD RELATIONS + INDEPENDENCE THREATS
	% Q355 -- Q384
	% ============================================================
	
	% ============================================================
	% Q355
	% ============================================================
	\begin{frame}[t]{Q355}
		\small
		\textbf{Question:} A CAE reports functionally to the audit committee and administratively to the CFO. The CFO is the subject of an anonymous hotline allegation of expense fraud. The CAE's only auditor with fraud experience is the CAE. The audit committee chair asks the CAE to ``handle it internally.'' Which is the MOST appropriate FIRST action?
		\begin{enumerate}[label=\Alph*.,itemsep=1pt,topsep=2pt]
			\item Conduct the investigation personally because the CAE is the only experienced fraud auditor.
			\item Recognize the self-review/independence conflict (CFO is administrative reporting line and subject), document it, and escalate to the audit committee with a recommendation for an independent investigator.
			\item Refer the matter to the external auditor and close the complaint.
			\item Ask the CFO for an explanation before deciding whether to investigate.
		\end{enumerate}
	\end{frame}
	
	\begin{frame}[t]{Q355 -- Answer \& Explanation}
		\small
		\textbf{Correct Answer: B} \\[4pt]
		\textbf{Why:} The CAE is administratively dependent on the CFO, and the CFO is the subject, creating a significant objectivity/independence threat. The FIRST step is to document and escalate to the audit committee with a recommendation for an independent investigator (Standards 1110/1130/1210.A2). \\[4pt]
		\textbf{A wrong:} The CAE investigating his own administrative superior is a self-review/conflict threat. \\[2pt]
		\textbf{C wrong:} Closing the complaint without investigation is improper. \\[2pt]
		\textbf{D wrong:} Asking the subject first tips off the investigation and compromises evidence.
	\end{frame}
	
	% ============================================================
	% Q356
	% ============================================================
	\begin{frame}[t]{Q356}
		\small
		\textbf{Question:} The CAE is asked by the CEO to sit on the credit committee that approves large loans. The CEO argues the CAE will ``have no vote, only observe,'' but the CAE would receive the same credit papers as voting members and be expected to comment on loan quality. The audit committee has not been consulted. Which is the MOST appropriate action?
		\begin{enumerate}[label=\Alph*.,itemsep=1pt,topsep=2pt]
			\item Accept, because a non-voting observer role cannot impair objectivity.
			\item Accept, because the CAE's presence strengthens the control environment and the audit committee can be informed later.
			\item Decline the voting/decision role but potentially accept a clearly advisory, non-decision role with audit committee awareness and safeguards against self-review of credit decisions later audited.
			\item Accept, but assign the future credit audit to a subordinate and document the arrangement.
		\end{enumerate}
	\end{frame}
	
	\begin{frame}[t]{Q356 -- Answer \& Explanation}
		\small
		\textbf{Correct Answer: C} \\[4pt]
		\textbf{Why:} Standards 1130 and 1130.A1/A2 prohibit assuming management responsibilities and require safeguards where advisory roles create self-review threats. Commenting on decisions and receiving decision papers can create self-review when those decisions are later audited. The audit committee must be aware. \\[4pt]
		\textbf{A wrong:} Non-voting does not eliminate self-review; commenting on loan quality is participation. \\[2pt]
		\textbf{B wrong:} Audit committee awareness cannot be deferred. \\[2pt]
		\textbf{D wrong:} Delegating the audit does not cure the CAE's self-review threat.
	\end{frame}
	
	% ============================================================
	% Q357
	% ============================================================
	\begin{frame}[t]{Q357}
		\small
		\textbf{Question:} The CAE's compensation is set by the CEO, who is also the sponsor of a major ERP program that internal audit is auditing. The CAE's proposed audit plan includes a review of the ERP program's governance. The CEO tells the CAE that including the ERP review ``will affect your bonus.'' Which is the MOST appropriate FIRST action?
		\begin{enumerate}[label=\Alph*.,itemsep=1pt,topsep=2pt]
			\item Remove the ERP review from the plan to protect the CAE's compensation.
			\item Document the threat and escalate to the audit committee, requesting that the committee address the compensation-linked threat and confirm the plan.
			\item Proceed with the audit but report findings only verbally to the CEO.
			\item Resign to avoid the conflict.
		\end{enumerate}
	\end{frame}
	
	\begin{frame}[t]{Q357 -- Answer \& Explanation}
		\small
		\textbf{Correct Answer: B} \\[4pt]
		\textbf{Why:} Linking audit coverage to the CAE's compensation is a direct intimidation/self-interest threat. The CAE must document and escalate to the audit committee, which should address the threat and confirm the plan (Standards 1110/1130/2010). \\[4pt]
		\textbf{A wrong:} Removing the review cedes the plan to CEO pressure. \\[2pt]
		\textbf{C wrong:} Verbal-only reporting to the subject of an audit is improper. \\[2pt]
		\textbf{D wrong:} Resignation is a last resort, not the first action.
	\end{frame}
	
	% ============================================================
	% Q358
	% ============================================================
	\begin{frame}[t]{Q358}
		\small
		\textbf{Question:} The audit committee chair privately tells the CAE that the CEO has ``verbally threatened'' to cut the internal audit budget if the plan includes a review of a high-margin business unit run by the CEO's son-in-law. The chair asks the CAE to ``handle it quietly.'' The unit represents 20\% of revenue and has never been audited. Which is the MOST appropriate action?
		\begin{enumerate}[label=\Alph*.,itemsep=1pt,topsep=2pt]
			\item Remove the unit from the plan to protect the budget and avoid conflict.
			\item Keep the unit in the plan based on risk, document the CEO's threat, and ensure the audit committee formally addresses the independence threat.
			\item Refer the matter to the external auditor and take no further action.
			\item Delay the audit for one year and revisit if the threat persists.
		\end{enumerate}
	\end{frame}
	
	\begin{frame}[t]{Q358 -- Answer \& Explanation}
		\small
		\textbf{Correct Answer: B} \\[4pt]
		\textbf{Why:} Standard 2010 requires a risk-based plan; Standards 1110/1111/1130 require the CAE to preserve independence and escalate threats. A CEO threat tied to coverage is a direct independence threat that must be documented and addressed formally with the audit committee. \\[4pt]
		\textbf{A wrong:} Removing a material, never-audited unit because of a threat is caving to intimidation. \\[2pt]
		\textbf{C wrong:} External audit referral does not discharge the CAE's responsibility. \\[2pt]
		\textbf{D wrong:} Delay rewards intimidation and leaves a material risk uncovered.
	\end{frame}
	
	% ============================================================
	% Q359
	% ============================================================
	\begin{frame}[t]{Q359}
		\small
		\textbf{Question:} A CAE has functional reporting to the audit committee but administrative reporting to the CFO. The CFO asks the CAE to ``hold off'' reporting a significant finance-department control deficiency until after the CFO's performance review. Which is the MOST appropriate action?
		\begin{enumerate}[label=\Alph*.,itemsep=1pt,topsep=2pt]
			\item Delay the report as requested to maintain a good working relationship.
			\item Report the deficiency promptly and escalate the CFO's request to the audit committee.
			\item Report the deficiency to the external auditor instead of the audit committee.
			\item Document the request and delay the report by one quarter.
		\end{enumerate}
	\end{frame}
	
	\begin{frame}[t]{Q359 -- Answer \& Explanation}
		\small
		\textbf{Correct Answer: B} \\[4pt]
		\textbf{Why:} Standard 2600 requires prompt reporting of significant issues to the board/audit committee. Administrative reporting to the CFO does not permit suppression, and the request itself is a governance concern. \\[4pt]
		\textbf{A wrong:} Delaying significant findings violates the Standards. \\[2pt]
		\textbf{C wrong:} Reporting to the external auditor does not replace board reporting. \\[2pt]
		\textbf{D wrong:} Documentation does not justify delay.
	\end{frame}
	
	% ============================================================
	% Q360
	% ============================================================
	\begin{frame}[t]{Q360}
		\small
		\textbf{Question:} The CAE is offered a paid speaking engagement by a vendor that internal audit is currently auditing. The fee is significant relative to the CAE's salary. Which is the MOST appropriate action?
		\begin{enumerate}[label=\Alph*.,itemsep=1pt,topsep=2pt]
			\item Accept; speaking fees are not gifts.
			\item Decline or disclose and evaluate whether it impairs objectivity, since it creates a self-interest threat.
			\item Accept but do not disclose.
			\item Accept only if the fee is small.
		\end{enumerate}
	\end{frame}
	
	\begin{frame}[t]{Q360 -- Answer \& Explanation}
		\small
		\textbf{Correct Answer: B} \\[4pt]
		\textbf{Why:} A paid engagement from an audited vendor creates a self-interest threat; decline or disclose and assess impairment (Code of Ethics, Objectivity). \\[4pt]
		\textbf{A wrong:} Fees can be a benefit with the same effect as gifts. \\[2pt]
		\textbf{C wrong:} Non-disclosure violates the Code. \\[2pt]
		\textbf{D wrong:} ``Small'' is subjective and still creates a threat.
	\end{frame}
	
	% ============================================================
	% Q361
	% ============================================================
	\begin{frame}[t]{Q361}
		\small
		\textbf{Question:} The CAE was the controller of the organization two years ago. The audit committee asks the CAE to lead an audit of the finance function, including controls the CAE designed as controller. Which is the MOST appropriate action?
		\begin{enumerate}[label=\Alph*.,itemsep=1pt,topsep=2pt]
			\item Accept, because two years have passed and the CAE now reports to the audit committee.
			\item Decline to lead the engagement and disclose the self-review threat; assign the engagement to an auditor without prior responsibility for the controls.
			\item Accept but use a smaller sample size to reduce the self-review risk.
			\item Accept and document the prior role in the workpapers.
		\end{enumerate}
	\end{frame}
	
	\begin{frame}[t]{Q361 -- Answer \& Explanation}
		\small
		\textbf{Correct Answer: B} \\[4pt]
		\textbf{Why:} Standard 1130.A1 prohibits auditing activities the auditor managed within the prior year; even beyond one year, designing the controls being audited creates a self-review threat. Decline to lead and reassign. \\[4pt]
		\textbf{A wrong:} Time and reporting line do not cure self-review. \\[2pt]
		\textbf{C wrong:} Smaller samples do not cure the threat. \\[2pt]
		\textbf{D wrong:} Documentation does not cure the threat.
	\end{frame}
	
	% ============================================================
	% Q362
	% ============================================================
	\begin{frame}[t]{Q362}
		\small
		\textbf{Question:} The CAE is asked by the CEO to lead the implementation of a new enterprise risk management (ERM) framework. The CEO argues ERM is a ``core internal audit competency'' and the CAE will ``own the framework.'' Which is the MOST appropriate response?
		\begin{enumerate}[label=\Alph*.,itemsep=1pt,topsep=2pt]
			\item Accept, because ERM is a core competency and the CEO supports it.
			\item Decline to own the framework; offer advisory/consulting support while management retains ownership, and avoid auditing the framework until objectivity is no longer impaired.
			\item Accept and later audit the framework as assurance.
			\item Decline any involvement because ERM is outside internal audit's scope.
		\end{enumerate}
	\end{frame}
	
	\begin{frame}[t]{Q362 -- Answer \& Explanation}
		\small
		\textbf{Correct Answer: B} \\[4pt]
		\textbf{Why:} Standards 1130 and 1130.A1/A2: internal audit may consult on ERM but must not own it or assume management responsibility. Owning it would impair objectivity and preclude later assurance. \\[4pt]
		\textbf{A wrong:} Competency is not the issue; ownership is. \\[2pt]
		\textbf{C wrong:} Self-review threat: implementing then auditing. \\[2pt]
		\textbf{D wrong:} Total abstention is overly restrictive; advisory support is permitted.
	\end{frame}
	
	% ============================================================
	% Q363
	% ============================================================
	\begin{frame}[t]{Q363}
		\small
		\textbf{Question:} The CAE is asked by the CFO to ``review the design'' of a new revenue system before go-live and to ``sign off that the design is adequate'' for the project steering committee. Which is the MOST appropriate response?
		\begin{enumerate}[label=\Alph*.,itemsep=1pt,topsep=2pt]
			\item Provide the review and sign off as requested, since the CFO is a key stakeholder.
			\item Provide advisory observations on design risks but decline to sign off, making clear that internal audit does not approve management's design or assume management responsibility.
			\item Decline any involvement because reviewing design impairs independence.
			\item Sign off but disclose in the report that the sign-off is advisory only.
		\end{enumerate}
	\end{frame}
	
	\begin{frame}[t]{Q363 -- Answer \& Explanation}
		\small
		\textbf{Correct Answer: B} \\[4pt]
		\textbf{Why:} Standards 1130 and the consulting guidance permit advisory review but prohibit assuming management responsibility or approving management's work. Sign-off implies approval/ownership of the design. \\[4pt]
		\textbf{A wrong:} Sign-off makes internal audit complicit and creates self-review. \\[2pt]
		\textbf{C wrong:} Advisory review is permitted; total abstention is not required. \\[2pt]
		\textbf{D wrong:} A disclaimer does not cure assumption of management responsibility.
	\end{frame}
	
	% ============================================================
	% Q364
	% ============================================================
	\begin{frame}[t]{Q364}
		\small
		\textbf{Question:} A CAE's spouse is the controller of a subsidiary that internal audit plans to audit. The CAE discloses the relationship to the audit committee chair, who says, ``As long as your spouse is not on the engagement team, it's fine.'' Which is the MOST accurate statement?
		\begin{enumerate}[label=\Alph*.,itemsep=1pt,topsep=2pt]
			\item The chair is correct; the CAE's spouse not being on the team eliminates the impairment.
			\item The impairment is not eliminated; a close family member in a key financial role at an audited subsidiary creates a familiarity/self-interest threat that requires reassigning the engagement or other safeguards approved by the full audit committee.
			\item The CAE should resign immediately.
			\item The CAE should proceed and document the relationship.
		\end{enumerate}
	\end{frame}
	
	\begin{frame}[t]{Q364 -- Answer \& Explanation}
		\small
		\textbf{Correct Answer: B} \\[4pt]
		\textbf{Why:} Standard 1130: objectivity is impaired when a close family member is in a position to influence the audited area. The chair alone cannot approve proceeding; safeguards (reassignment, independent oversight) are required and the full committee should be involved. \\[4pt]
		\textbf{A wrong:} Not being on the team is necessary but not sufficient. \\[2pt]
		\textbf{C wrong:} Resignation is an overreaction; safeguards are available. \\[2pt]
		\textbf{D wrong:} Documentation alone does not cure the threat.
	\end{frame}
	
	% ============================================================
	% Q365
	% ============================================================
	\begin{frame}[t]{Q365}
		\small
		\textbf{Question:} The CAE is asked to approve the audit committee's proposed change to the internal audit charter that would require the CAE to obtain CEO approval before communicating results to the audit committee. Which is the MOST appropriate response?
		\begin{enumerate}[label=\Alph*.,itemsep=1pt,topsep=2pt]
			\item Accept the change because the audit committee proposed it.
			\item Advise that the change would impair the CAE's ability to communicate freely with the audit committee and contradict Standard 1111/2600; the CAE should explain the impairment and request the committee reconsider.
			\item Accept but communicate results to the CEO only.
			\item Resign.
		\end{enumerate}
	\end{frame}
	
	\begin{frame}[t]{Q365 -- Answer \& Explanation}
		\small
		\textbf{Correct Answer: B} \\[4pt]
		\textbf{Why:} Direct, unrestricted communication with the audit committee is essential to independence (Standards 1111/2600). A CEO gatekeeping requirement would impair that. The CAE must advise against it and explain the impairment. \\[4pt]
		\textbf{A wrong:} Accepting an impairing charter change is improper. \\[2pt]
		\textbf{C wrong:} Communicating only to the CEO further impairs independence. \\[2pt]
		\textbf{D wrong:} Resignation is a last resort.
	\end{frame}
	
	% ============================================================
	% Q366
	% ============================================================
	\begin{frame}[t]{Q366}
		\small
		\textbf{Question:} The audit committee asks the CAE to evaluate the CEO's expense reimbursements. The CEO approves the CAE's budget and has previously threatened to ``restructure'' internal audit if findings about executives were reported. The audit committee chair assures the CAE that ``the committee will back you.'' Which is the MOST appropriate action?
		\begin{enumerate}[label=\Alph*.,itemsep=1pt,topsep=2pt]
			\item Decline because the CEO's budget control and threat impair independence.
			\item Accept, document the threat, ensure the audit committee receives results directly, and consider assigning a strong engagement team to reduce dependence on the CEO.
			\item Accept but report findings only verbally to the chair.
			\item Accept and exclude the CEO's travel expenses from scope.
		\end{enumerate}
	\end{frame}
	
	\begin{frame}[t]{Q366 -- Answer \& Explanation}
		\small
		\textbf{Correct Answer: B} \\[4pt]
		\textbf{Why:} The audit committee provides functional oversight and has committed support; safeguards (direct reporting, documentation, competent team) mitigate the threat. Standards 1110/1130/2600 support accepting with safeguards rather than declining. \\[4pt]
		\textbf{A wrong:} Declining leaves a governance-requested engagement undone; safeguards can mitigate. \\[2pt]
		\textbf{C wrong:} Verbal-only reporting is not an appropriate communication channel. \\[2pt]
		\textbf{D wrong:} Excluding the CEO's travel from scope defeats the engagement.
	\end{frame}
	
	% ============================================================
	% Q367
	% ============================================================
	\begin{frame}[t]{Q367}
		\small
		\textbf{Question:} The CAE reports functionally to the audit committee and administratively to the CFO. The CFO asks the CAE to perform a ``quick review'' of a peer company's financial statements to help the CFO's personal investment decision. Which is the MOST appropriate response?
		\begin{enumerate}[label=\Alph*.,itemsep=1pt,topsep=2pt]
			\item Accept, because the CFO is the administrative reporting line.
			\item Decline, because the request is for personal benefit, outside the charter, and creates an objectivity/independence threat; inform the audit committee if pressure persists.
			\item Accept but do not document the engagement.
			\item Accept only if the review is verbal.
		\end{enumerate}
	\end{frame}
	
	\begin{frame}[t]{Q367 -- Answer \& Explanation}
		\small
		\textbf{Correct Answer: B} \\[4pt]
		\textbf{Why:} Using internal audit resources for the CFO's personal investment decision is outside the charter, for personal benefit, and creates an objectivity threat. Decline; escalate if pressure persists. \\[4pt]
		\textbf{A wrong:} Administrative reporting does not authorize personal use of internal audit. \\[2pt]
		\textbf{C wrong:} Undocumented engagements violate the Standards. \\[2pt]
		\textbf{D wrong:} Verbal delivery does not cure the threat.
	\end{frame}
	
	% ============================================================
	% Q368
	% ============================================================
	\begin{frame}[t]{Q368}
		\small
		\textbf{Question:} The CAE is evaluating a request from the audit committee chair to ``co-sign'' an engagement report prepared by an external consultant who performed an internal audit outsourcing engagement. The consultant's report concluded ``no significant issues.'' The CAE did not supervise or review the consultant's work. Which is the MOST appropriate response?
		\begin{enumerate}[label=\Alph*.,itemsep=1pt,topsep=2pt]
			\item Co-sign the report because the chair requested it.
			\item Decline to co-sign unless the CAE has sufficient evidence of the consultant's work quality and can take responsibility for the conclusions; otherwise, provide a separate statement of the CAE's role.
			\item Co-sign but add a disclaimer.
			\item Co-sign only if the consultant is a CPA firm.
		\end{enumerate}
	\end{frame}
	
	\begin{frame}[t]{Q368 -- Answer \& Explanation}
		\small
		\textbf{Correct Answer: B} \\[4pt]
		\textbf{Why:} Signing a report implies responsibility for its content (Standards 2330/2410). The CAE cannot co-sign without evidence of work quality and the ability to take responsibility. \\[4pt]
		\textbf{A wrong:} Chair request does not justify signing for unperformed work. \\[2pt]
		\textbf{C wrong:} A disclaimer does not cure the lack of basis. \\[2pt]
		\textbf{D wrong:} CPA status is not the criterion.
	\end{frame}
	
	% ============================================================
	% Q369
	% ============================================================
	\begin{frame}[t]{Q369}
		\small
		\textbf{Question:} The CAE is a member of the organization's risk committee, which approves risk responses and risk appetite recommendations to the board. The CAE has voting rights on the committee. Which is the MOST accurate statement?
		\begin{enumerate}[label=\Alph*.,itemsep=1pt,topsep=2pt]
			\item Voting membership is appropriate because it strengthens internal audit's influence.
			\item Voting membership impairs objectivity because the CAE would be approving the risk responses that internal audit later assesses; non-voting/advisory membership with disclosure is preferable.
			\item Voting membership is required by the Standards.
			\item Voting membership is acceptable if the CAE discloses it in the annual report.
		\end{enumerate}
	\end{frame}
	
	\begin{frame}[t]{Q369 -- Answer \& Explanation}
		\small
		\textbf{Correct Answer: B} \\[4pt]
		\textbf{Why:} Voting on risk responses the CAE will later audit creates a self-review/management-responsibility threat (Standards 1130/1130.A1). Non-voting/advisory membership is preferable. \\[4pt]
		\textbf{A wrong:} Influence does not justify impaired objectivity. \\[2pt]
		\textbf{C wrong:} The Standards do not require voting membership. \\[2pt]
		\textbf{D wrong:} Disclosure does not cure the self-review threat.
	\end{frame}
	
	% ============================================================
	% Q370
	% ============================================================
	\begin{frame}[t]{Q370}
		\small
		\textbf{Question:} A new CEO tells the CAE that, going forward, all audit reports must be routed through the CEO's office before being sent to the audit committee, so the CEO can ``add context.'' The charter currently gives the CAE direct access to the audit committee. Which is the MOST appropriate action?
		\begin{enumerate}[label=\Alph*.,itemsep=1pt,topsep=2pt]
			\item Comply, because the CEO is the administrative reporting line.
			\item Decline to cede direct access; explain that the charter and Standards 1111/2600 require direct communication with the audit committee, and escalate the matter to the audit committee.
			\item Comply but send a copy of reports to the audit committee in parallel.
			\item Resign.
		\end{enumerate}
	\end{frame}
	
	\begin{frame}[t]{Q370 -- Answer \& Explanation}
		\small
		\textbf{Correct Answer: B} \\[4pt]
		\textbf{Why:} Direct access to the audit committee is essential to independence (Standards 1111/2600). CEO gatekeeping impairs it; the CAE must decline and escalate. \\[4pt]
		\textbf{A wrong:} Administrative reporting does not authorize gatekeeping. \\[2pt]
		\textbf{C wrong:} Parallel copies may still be filtered/modified; direct access is required. \\[2pt]
		\textbf{D wrong:} Resignation is a last resort.
	\end{frame}
	
	% ============================================================
	% Q371
	% ============================================================
	\begin{frame}[t]{Q371}
		\small
		\textbf{Question:} The audit committee asks the CAE to perform a consulting engagement to redesign the organization's procurement approval workflow. The CAE accepts and completes the redesign. One year later, the audit committee asks internal audit to audit the procurement approval workflow. Which is the MOST appropriate action?
		\begin{enumerate}[label=\Alph*.,itemsep=1pt,topsep=2pt]
			\item Accept the audit; one year has passed since the redesign.
			\item Decline to lead the audit and disclose the self-review threat; assign the audit to an auditor without prior responsibility for the redesign, with independent oversight.
			\item Accept but use a smaller sample.
			\item Accept and document the redesign in the workpapers.
		\end{enumerate}
	\end{frame}
	
	\begin{frame}[t]{Q371 -- Answer \& Explanation}
		\small
		\textbf{Correct Answer: B} \\[4pt]
		\textbf{Why:} Designing the workflow creates a self-review threat for auditing it (Standards 1130/1130.A1). Reassignment with independent oversight is required; time passage alone does not cure the threat. \\[4pt]
		\textbf{A wrong:} One year does not automatically cure self-review. \\[2pt]
		\textbf{C wrong:} Sample size does not cure the threat. \\[2pt]
		\textbf{D wrong:} Documentation does not cure the threat.
	\end{frame}
	
	% ============================================================
	% Q372
	% ============================================================
	\begin{frame}[t]{Q372}
		\small
		\textbf{Question:} The CAE learns that the audit committee chair's brother is the CFO of a major vendor. The audit committee chair participates in discussions about internal audit's coverage of vendor management. Which is the MOST appropriate action?
		\begin{enumerate}[label=\Alph*.,itemsep=1pt,topsep=2pt]
			\item Ignore it because the chair is not an employee.
			\item Disclose the relationship to the full audit committee and evaluate whether the chair should recuse from vendor-management oversight; document the consideration.
			\item Confront the chair publicly at the next meeting.
			\item Report the matter to the external auditor only.
		\end{enumerate}
	\end{frame}
	
	\begin{frame}[t]{Q372 -- Answer \& Explanation}
		\small
		\textbf{Correct Answer: B} \\[4pt]
		\textbf{Why:} The chair's family relationship with a major vendor creates a conflict of interest affecting oversight of vendor management. Disclosure and recusal evaluation are required. \\[4pt]
		\textbf{A wrong:} Board members are subject to conflict-of-interest considerations. \\[2pt]
		\textbf{C wrong:} Public confrontation is unprofessional. \\[2pt]
		\textbf{D wrong:} External auditor is not the appropriate channel.
	\end{frame}
	
	% ============================================================
	% Q373
	% ============================================================
	\begin{frame}[t]{Q373}
		\small
		\textbf{Question:} The CAE is asked by the audit committee to assess whether the organization's whistleblower hotline is effective. The hotline is operated by an external provider and reports go to the audit committee chair, who decides whether to forward them to internal audit. Which is the MOST significant weakness?
		\begin{enumerate}[label=\Alph*.,itemsep=1pt,topsep=2pt]
			\item The hotline is operated by an external provider.
			\item Reports are filtered through the audit committee chair, who may suppress matters and may be the subject of a complaint; the full committee and internal audit should have independent access.
			\item The hotline is not operated by internal audit.
			\item The hotline is not operated by management.
		\end{enumerate}
	\end{frame}
	
	\begin{frame}[t]{Q373 -- Answer \& Explanation}
		\small
		\textbf{Correct Answer: B} \\[4pt]
		\textbf{Why:} Filtering through a single person (especially the chair, who could be a subject) creates a suppression risk. Independent access by the full committee and internal audit is required for effectiveness. \\[4pt]
		\textbf{A wrong:} External operation is a strength. \\[2pt]
		\textbf{C wrong:} Internal audit operation is not required. \\[2pt]
		\textbf{D wrong:} Management operation is not required.
	\end{frame}
	
	% ============================================================
	% Q374
	% ============================================================
	\begin{frame}[t]{Q374}
		\small
		\textbf{Question:} The CAE is asked to provide assurance on the audit committee's own effectiveness. The CAE has direct access to the committee, and the committee approves the CAE's charter, plan, and compensation. Which is the MOST appropriate action?
		\begin{enumerate}[label=\Alph*.,itemsep=1pt,topsep=2pt]
			\item Decline the engagement because the CAE reports to the committee.
			\item Accept the engagement with safeguards: use an independent reviewer or external facilitator, disclose the self-review threat, and ensure the methodology is objective.
			\item Accept the engagement and perform it alone.
			\item Report the committee's effectiveness only verbally.
		\end{enumerate}
	\end{frame}
	
	\begin{frame}[t]{Q374 -- Answer \& Explanation}
		\small
		\textbf{Correct Answer: B} \\[4pt]
		\textbf{Why:} Assessing the committee that oversees the CAE creates a self-review/independence threat. Safeguards (independent reviewer, disclosure, objective methodology) are required; the engagement is not prohibited. \\[4pt]
		\textbf{A wrong:} Declining is unnecessary; safeguards are available. \\[2pt]
		\textbf{C wrong:} Performing it alone without safeguards is improper. \\[2pt]
		\textbf{D wrong:} Verbal-only reporting is not adequate.
	\end{frame}
	
	% ============================================================
	% Q375
	% ============================================================
	\begin{frame}[t]{Q375}
		\small
		\textbf{Question:} The CAE reports functionally to the audit committee and administratively to the CFO. The CFO asks the CAE to ``review my department's budget request to ensure it is defensible'' before submission to the board. Which is the MOST appropriate response?
		\begin{enumerate}[label=\Alph*.,itemsep=1pt,topsep=2pt]
			\item Accept, because the CFO is the administrative reporting line.
			\item Decline to endorse the CFO's budget; explain that internal audit can provide independent observations on the budget process but cannot advocate for management's position, which would impair objectivity.
			\item Accept but disclose the request in the next report.
			\item Accept only if the review is verbal.
		\end{enumerate}
	\end{frame}
	
	\begin{frame}[t]{Q375 -- Answer \& Explanation}
		\small
		\textbf{Correct Answer: B} \\[4pt]
		\textbf{Why:} Advocating for management's budget position compromises objectivity (Code, Standards 1130). Internal audit may observe the budget process independently but not endorse management's submission. \\[4pt]
		\textbf{A wrong:} Administrative reporting does not authorize advocacy. \\[2pt]
		\textbf{C wrong:} Disclosure after the fact does not cure the impairment. \\[2pt]
		\textbf{D wrong:} Verbal delivery does not cure the impairment.
	\end{frame}
	
	% ============================================================
	% Q376
	% ============================================================
	\begin{frame}[t]{Q376}
		\small
		\textbf{Question:} The audit committee asks the CAE to investigate an allegation that the CEO violated the code of conduct. The CEO is the CAE's administrative reporting line and has a close personal relationship with the audit committee chair. Which is the MOST appropriate action?
		\begin{enumerate}[label=\Alph*.,itemsep=1pt,topsep=2pt]
			\item Accept the investigation and report only to the chair.
			\item Accept the engagement with safeguards: report to the full audit committee (not the chair alone), document independence considerations, and consider an independent investigator if the chair's relationship impairs oversight.
			\item Decline because the CEO is the administrative reporting line.
			\item Refer the matter to the external auditor and close the matter.
		\end{enumerate}
	\end{frame}
	
	\begin{frame}[t]{Q376 -- Answer \& Explanation}
		\small
		\textbf{Correct Answer: B} \\[4pt]
		\textbf{Why:} The chair's close relationship with the CEO impairs the chair's oversight, so reporting to the full committee (and considering an independent investigator) is required. Accepting with safeguards is appropriate. \\[4pt]
		\textbf{A wrong:} Reporting only to the conflicted chair is improper. \\[2pt]
		\textbf{C wrong:} Declining leaves a governance matter unaddressed; safeguards are available. \\[2pt]
		\textbf{D wrong:} External referral does not discharge internal audit's responsibility.
	\end{frame}
	
% ============================================================
% Q377
% ============================================================
\begin{frame}[t]{Q377}
	\small
	\textbf{Question:} The CAE is asked by the CEO to ``review the external auditor's workpapers'' to help the CEO negotiate a lower audit fee. The external auditor has not authorized access. Which is the MOST appropriate action?
	\begin{enumerate}[label=\Alph*.,itemsep=1pt,topsep=2pt]
		\item Accept, because the CEO has authority over vendor relationships.
		\item Decline, because the external auditor's workpapers are confidential without authorization, and using them for fee negotiation is outside internal audit's scope; the audit committee should be informed if pressure persists.
		\item Accept and share the workpapers with the CEO.
		\item Accept but do not document the review.
	\end{enumerate}
\end{frame}

\begin{frame}[t]{Q377 -- Answer \& Explanation}
	\small
	\textbf{Correct Answer: B} \\[4pt]
	\textbf{Why:} External auditor workpapers are confidential without authorization. Using them for fee negotiation is outside internal audit's scope and creates an ethical/confidentiality issue. Decline and escalate if pressured. \\[4pt]
	\textbf{A wrong:} CEO authority does not override confidentiality. \\[2pt]
	\textbf{C wrong:} Sharing without authorization breaches confidentiality. \\[2pt]
	\textbf{D wrong:} Undocumented engagements violate the Standards.
\end{frame}
	
	% ============================================================
	% Q378
	% ============================================================
	\begin{frame}[t]{Q378}
		\small
		\textbf{Question:} The CAE is a member of the organization's ethics committee, which adjudicates ethics violations and imposes discipline. Which is the MOST accurate statement?
		\begin{enumerate}[label=\Alph*.,itemsep=1pt,topsep=2pt]
			\item Membership is appropriate because it strengthens ethics oversight.
			\item Membership may impair objectivity if the CAE would later audit ethics investigations or discipline decisions; an advisory role with disclosure is preferable.
			\item Membership is required by the Standards.
			\item Membership is acceptable if the CAE is not the committee chair.
		\end{enumerate}
	\end{frame}
	
	\begin{frame}[t]{Q378 -- Answer \& Explanation}
		\small
		\textbf{Correct Answer: B} \\[4pt]
		\textbf{Why:} Adjudicating discipline the CAE may later audit creates a self-review/management-responsibility threat (Standards 1130/1130.A1). Advisory membership with disclosure is preferable. \\[4pt]
		\textbf{A wrong:} Influence does not justify impaired objectivity. \\[2pt]
		\textbf{C wrong:} The Standards do not require membership. \\[2pt]
		\textbf{D wrong:} Not being chair does not cure the threat.
	\end{frame}
	
	% ============================================================
	% Q379
	% ============================================================
	\begin{frame}[t]{Q379}
		\small
		\textbf{Question:} The CAE is asked by the audit committee to evaluate the internal audit function's compliance with the Standards. The CAE finds that the only nonconformance is that the external quality assessment was performed 6 years ago. Which is the MOST appropriate statement?
		\begin{enumerate}[label=\Alph*.,itemsep=1pt,topsep=2pt]
			\item The activity ``generally conforms'' because the only issue is the external assessment.
			\item The activity does not fully conform to Standard 1312; the CAE must disclose the nonconformance to the audit committee and develop a remediation plan (Standard 1320).
			\item The external assessment is optional if internal assessments are clean.
			\item The CAE should ignore the issue because it is only one nonconformance.
		\end{enumerate}
	\end{frame}
	
	\begin{frame}[t]{Q379 -- Answer \& Explanation}
		\small
		\textbf{Correct Answer: B} \\[4pt]
		\textbf{Why:} Standard 1312 requires external assessment at least every five years. A 6-year gap is a nonconformance requiring disclosure to the board and a remediation plan (Standard 1320). \\[4pt]
		\textbf{A wrong:} A required element is missing; ``generally conforms'' is not justified. \\[2pt]
		\textbf{C wrong:} External assessment is mandatory. \\[2pt]
		\textbf{D wrong:} Ignoring a nonconformance is improper.
	\end{frame}
	
	% ============================================================
	% Q380
	% ============================================================
	\begin{frame}[t]{Q380}
		\small
		\textbf{Question:} The CAE is evaluating whether to accept an engagement requested by the CEO to review the CEO's own expense reports. The audit committee has not been informed, and the CEO states the engagement is ``confidential.'' Which is the MOST appropriate action?
		\begin{enumerate}[label=\Alph*.,itemsep=1pt,topsep=2pt]
			\item Accept because the CEO has authority to direct internal audit.
			\item Decline unless the audit committee is informed and functional oversight is provided; auditing the CEO's expenses requires committee awareness and safeguards.
			\item Accept and report results only to the CEO.
			\item Accept and assign the engagement to the CAE personally.
		\end{enumerate}
	\end{frame}
	
	\begin{frame}[t]{Q380 -- Answer \& Explanation}
		\small
		\textbf{Correct Answer: B} \\[4pt]
		\textbf{Why:} Auditing the CEO's expenses is a sensitive engagement requiring audit committee awareness and functional oversight (Standards 1110/1130/2600). CEO-confidential handling is improper. \\[4pt]
		\textbf{A wrong:} CEO authority does not override audit committee oversight. \\[2pt]
		\textbf{C wrong:} Reporting only to the subject is improper. \\[2pt]
		\textbf{D wrong:} CAE personally performing it does not cure the oversight gap.
	\end{frame}
	
	% ============================================================
	% Q381
	% ============================================================
	\begin{frame}[t]{Q381}
		\small
		\textbf{Question:} The CAE is asked by the audit committee to evaluate the organization's compliance with a new data privacy law. The CAE has no privacy expertise, and the audit committee wants the engagement completed in 30 days. Which is the MOST appropriate action?
		\begin{enumerate}[label=\Alph*.,itemsep=1pt,topsep=2pt]
			\item Accept and perform the engagement personally.
			\item Accept with a realistic plan that includes obtaining competent privacy expertise (internal specialist or external consultant), and communicate the resource/timing implications to the audit committee.
			\item Decline because internal audit lacks privacy expertise.
			\item Accept and reduce the scope to areas the CAE understands.
		\end{enumerate}
	\end{frame}
	
	\begin{frame}[t]{Q381 -- Answer \& Explanation}
		\small
		\textbf{Correct Answer: B} \\[4pt]
		\textbf{Why:} Standard 1210 permits obtaining competent advice/specialists. The CAE should accept with a realistic plan, obtain expertise, and communicate resource/timing implications (Standard 2030). \\[4pt]
		\textbf{A wrong:} Performing without competence violates the Standards. \\[2pt]
		\textbf{C wrong:} Declining is unnecessary; expertise can be obtained. \\[2pt]
		\textbf{D wrong:} Narrowing scope to fit competence defeats the engagement.
	\end{frame}
	
	% ============================================================
	% Q382
	% ============================================================
	\begin{frame}[t]{Q382}
		\small
		\textbf{Question:} A board member asks the CAE to investigate a personal business dispute involving the board member's family. The board member offers to ``make it worth your while.'' Which is the MOST appropriate action?
		\begin{enumerate}[label=\Alph*.,itemsep=1pt,topsep=2pt]
			\item Accept because board members are governance officials.
			\item Decline, because the request is for personal benefit, outside the charter, and the offer creates an integrity/objectivity threat; report the offer to the audit committee.
			\item Accept but do not document the engagement.
			\item Accept only if the investigation is brief.
		\end{enumerate}
	\end{frame}
	
	\begin{frame}[t]{Q382 -- Answer \& Explanation}
		\small
		\textbf{Correct Answer: B} \\[4pt]
		\textbf{Why:} The request is personal, outside the charter, and the offer is a bribe attempt. Decline and report the offer (Code Integrity/Objectivity; Standards 1000/1130). \\[4pt]
		\textbf{A wrong:} Board status does not authorize personal use of internal audit. \\[2pt]
		\textbf{C wrong:} Undocumented engagements violate the Standards. \\[2pt]
		\textbf{D wrong:} Duration does not cure the impropriety.
	\end{frame}
	
	% ============================================================
	% Q383
	% ============================================================
	\begin{frame}[t]{Q383}
		\small
		\textbf{Question:} The CAE is asked by the audit committee to assess the board's effectiveness in overseeing risk. The CAE has direct access to the board but not to board self-assessment records. Which is the MOST appropriate action?
		\begin{enumerate}[label=\Alph*.,itemsep=1pt,topsep=2pt]
			\item Decline because the CAE lacks access to board records.
			\item Accept with a methodology that combines available evidence (board minutes, committee reports, interviews with directors) and clearly state scope limitations in the report.
			\item Accept and conclude the board is effective based on interviews alone.
			\item Report scope limitations only if challenged.
		\end{enumerate}
	\end{frame}
	
	\begin{frame}[t]{Q383 -- Answer \& Explanation}
		\small
		\textbf{Correct Answer: B} \\[4pt]
		\textbf{Why:} The CAE can assess board risk oversight using available evidence and must clearly communicate scope limitations (Standards 2200/2410). Declining is unnecessary; interviews alone are insufficient. \\[4pt]
		\textbf{A wrong:} The engagement is feasible with available evidence. \\[2pt]
		\textbf{C wrong:} Interviews alone are insufficient. \\[2pt]
		\textbf{D wrong:} Scope limitations must be communicated in the report.
	\end{frame}
	
	% ============================================================
	% Q384
	% ============================================================
	\begin{frame}[t]{Q384}
		\small
		\textbf{Question:} The CAE's performance evaluation and compensation are approved by the CFO, while the audit committee only ``provides input.'' Which is the MOST accurate statement?
		\begin{enumerate}[label=\Alph*.,itemsep=1pt,topsep=2pt]
			\item This arrangement is appropriate because the audit committee provides input.
			\item This arrangement impairs independence; Standard 1111 requires the audit committee to approve the CAE's performance evaluation and compensation (functional reporting).
			\item This arrangement is required by the Standards.
			\item This arrangement is acceptable if the CAE is a CIA.
		\end{enumerate}
	\end{frame}
	
	\begin{frame}[t]{Q384 -- Answer \& Explanation}
		\small
		\textbf{Correct Answer: B} \\[4pt]
		\textbf{Why:} Standard 1111: functional reporting to the board/audit committee includes approving the CAE's performance evaluation and compensation. CFO approval with committee ``input'' impairs independence. \\[4pt]
		\textbf{A wrong:} Input is not approval. \\[2pt]
		\textbf{C wrong:} The Standards require committee approval. \\[2pt]
		\textbf{D wrong:} Certification does not cure the structural impairment.
	\end{frame}
	
	% ============================================================
	% DOMAIN B: QAIP + CONFORMANCE JUDGMENT + EXTERNAL ASSESSMENT
	% Q385 -- Q413
	% ============================================================
	
	% ============================================================
	% Q385
	% ============================================================
	\begin{frame}[t]{Q385}
		\small
		\textbf{Question:} The CAE is assessing the internal audit activity's conformance. The activity has: (1) a board-approved charter; (2) a risk-based plan; (3) no QAIP; (4) an engagement that audited a process the engagement partner designed; (5) a CAE who reports functionally to the audit committee. The CAE wants to conclude ``generally conforms.'' Which issue MOST directly precludes that conclusion?
		\begin{enumerate}[label=\Alph*.,itemsep=1pt,topsep=2pt]
			\item The absence of a QAIP.
			\item The self-review engagement (auditing a process the partner designed).
			\item The risk-based plan.
			\item The functional reporting to the audit committee.
		\end{enumerate}
	\end{frame}
	
	\begin{frame}[t]{Q385 -- Answer \& Explanation}
		\small
		\textbf{Correct Answer: A} \\[4pt]
		\textbf{Why:} Both the absence of a QAIP (Standards 1300–1312) and the self-review engagement (Standard 1130.A1) are nonconformances. The QAIP gap is systemic (activity-wide); the self-review is engagement-specific. A systemic mandatory requirement missing precludes ``generally conforms.'' \\[4pt]
		\textbf{B wrong:} Serious, but engagement-specific. \\[2pt]
		\textbf{C wrong:} Risk-based plan is a strength. \\[2pt]
		\textbf{D wrong:} Functional reporting is a strength.
	\end{frame}
	
	% ============================================================
	% Q386
	% ============================================================
	\begin{frame}[t]{Q386}
		\small
		\textbf{Question:} The internal audit activity has a QAIP that includes ongoing internal assessments but no periodic internal assessments, and the external assessment was completed 5 years and 2 months ago. Which is the MOST accurate statement?
		\begin{enumerate}[label=\Alph*.,itemsep=1pt,topsep=2pt]
			\item The QAIP conforms because ongoing assessments exist and the external assessment is ``about'' every five years.
			\item The QAIP does not fully conform; periodic internal assessments are required (Standard 1311) and the external assessment is beyond the five-year requirement (Standard 1312).
			\item Only the external assessment gap matters.
			\item The QAIP conforms because no significant issues were found.
		\end{enumerate}
	\end{frame}
	
	\begin{frame}[t]{Q386 -- Answer \& Explanation}
		\small
		\textbf{Correct Answer: B} \\[4pt]
		\textbf{Why:} Standard 1311 requires both ongoing and periodic internal assessments; Standard 1312 requires external assessment at least every five years. Missing periodic assessments and exceeding the five-year window are both nonconformances. \\[4pt]
		\textbf{A wrong:} ``About'' five years does not satisfy ``at least every five years.'' \\[2pt]
		\textbf{C wrong:} Both gaps matter. \\[2pt]
		\textbf{D wrong:} Absence of findings does not waive requirements.
	\end{frame}
	
	% ============================================================
	% Q387
	% ============================================================
	\begin{frame}[t]{Q387}
		\small
		\textbf{Question:} The CAE is selecting an external quality assessor. The candidate is a former partner of the CAE's spouse, is highly qualified, and is independent of the organization. Which is the MOST appropriate action?
		\begin{enumerate}[label=\Alph*.,itemsep=1pt,topsep=2pt]
			\item Engage the candidate because qualifications and independence from the organization are met.
			\item Do not engage the candidate without disclosure and evaluation of the relationship; the former partner of the CAE's spouse creates an independence threat for the assessment.
			\item Engage the candidate and disclose after the assessment.
			\item Engage the candidate only if the audit committee approves.
		\end{enumerate}
	\end{frame}
	
	\begin{frame}[t]{Q387 -- Answer \& Explanation}
		\small
		\textbf{Correct Answer: B} \\[4pt]
		\textbf{Why:} Standards 1312 and the QAIP require the external assessor to be independent and objective. A former partner of the CAE's spouse is a significant independence threat; disclosure and evaluation (or selecting another assessor) are required. \\[4pt]
		\textbf{A wrong:} Independence from the organization does not equal independence from the CAE. \\[2pt]
		\textbf{C wrong:} Post-hoc disclosure does not cure the impairment. \\[2pt]
		\textbf{D wrong:} Committee approval does not cure the independence threat.
	\end{frame}
	
	% ============================================================
	% Q388
	% ============================================================
	\begin{frame}[t]{Q388}
		\small
		\textbf{Question:} The external quality assessment concludes ``does not conform'' because the CAE reports administratively to the CFO and the audit committee has no authority over the CAE's appointment. Which is the MOST appropriate action?
		\begin{enumerate}[label=\Alph*.,itemsep=1pt,topsep=2pt]
			\item Ignore the finding because the CAE's day-to-day work is independent.
			\item Develop a remediation plan (e.g., charter change to give the audit committee authority over the CAE's appointment/evaluation), report the nonconformance and plan to the board, and monitor implementation (Standard 1320).
			\item Dispute the finding because administrative reporting to the CFO is common.
			\item Resign.
		\end{enumerate}
	\end{frame}
	
	\begin{frame}[t]{Q388 -- Answer \& Explanation}
		\small
		\textbf{Correct Answer: B} \\[4pt]
		\textbf{Why:} Standard 1320 requires reporting significant nonconformance to senior management and the board and developing a remediation plan. The structural independence issue must be addressed (charter/committee authority). \\[4pt]
		\textbf{A wrong:} Ignoring the finding is improper. \\[2pt]
		\textbf{C wrong:} Common practice does not cure a nonconformance. \\[2pt]
		\textbf{D wrong:} Resignation is a last resort.
	\end{frame}
	
	% ============================================================
	% Q389
	% ============================================================
	\begin{frame}[t]{Q389}
		\small
		\textbf{Question:} The CAE is evaluating whether the internal audit activity can conclude ``conforms'' with the Standards. The activity has a QAIP, ongoing and periodic internal assessments, and an external assessment within five years. However, the CAE has not disclosed a significant nonconformance (an engagement that audited the CAE's prior department within six months) to the board. Which is the MOST accurate statement?
		\begin{enumerate}[label=\Alph*.,itemsep=1pt,topsep=2pt]
			\item The activity conforms because the QAIP structure is complete.
			\item The activity does not conform; the undisclosed self-review nonconformance and the failure to report it to the board violate Standards 1130.A1 and 1320.
			\item The self-review is acceptable because it was disclosed internally to the engagement team.
			\item The external assessment cures the nonconformance.
		\end{enumerate}
	\end{frame}
	
	\begin{frame}[t]{Q389 -- Answer \& Explanation}
		\small
		\textbf{Correct Answer: B} \\[4pt]
		\textbf{Why:} Standard 1130.A1 prohibits self-review (auditing the prior department within one year); Standard 1320 requires reporting nonconformance to the board. Structure alone does not cure substance. \\[4pt]
		\textbf{A wrong:} Structure does not cure the self-review and non-disclosure. \\[2pt]
		\textbf{C wrong:} Internal team disclosure is not board reporting. \\[2pt]
		\textbf{D wrong:} External assessment does not cure nonconformance.
	\end{frame}
	
	% ============================================================
	% Q390
	% ============================================================
	\begin{frame}[t]{Q390}
		\small
		\textbf{Question:} The CAE is evaluating the QAIP's ongoing internal assessments. The activity uses engagement-level supervision, a checklist completed by engagement leads, and an annual survey of auditees. No independent post-engagement review is performed by someone not on the engagement team. Which is the MOST accurate statement?
		\begin{enumerate}[label=\Alph*.,itemsep=1pt,topsep=2pt]
			\item The ongoing assessments are sufficient because supervision and checklists exist.
			\item The ongoing assessments are weak because independent post-engagement review by someone not on the team is a key element of ongoing internal assessments (Standard 1311).
			\item Auditee surveys alone satisfy ongoing assessment requirements.
			\item Engagement lead checklists alone satisfy ongoing assessment requirements.
		\end{enumerate}
	\end{frame}
	
	\begin{frame}[t]{Q390 -- Answer \& Explanation}
		\small
		\textbf{Correct Answer: B} \\[4pt]
		\textbf{Why:} Standard 1311 ongoing assessments should include independent review of engagement work by someone not on the team (e.g., engagement-level supervision and review, checklists, and post-engagement reviews). Self-completed checklists and surveys alone are insufficient. \\[4pt]
		\textbf{A wrong:} Supervision and checklists are necessary but not sufficient without independent review. \\[2pt]
		\textbf{C wrong:} Surveys alone are insufficient. \\[2pt]
		\textbf{D wrong:} Self-checklists alone are insufficient.
	\end{frame}
	
	% ============================================================
	% Q391
	% ============================================================
	\begin{frame}[t]{Q391}
		\small
		\textbf{Question:} The external quality assessor is a firm that also provides consulting services to the organization's internal audit function (e.g., training and methodology design). Which is the MOST significant concern?
		\begin{enumerate}[label=\Alph*.,itemsep=1pt,topsep=2pt]
			\item The firm's qualifications.
			\item The firm's independence/objectivity, because providing consulting services to the function being assessed creates a self-review threat.
			\item The cost of the assessment.
			\item The duration of the assessment.
		\end{enumerate}
	\end{frame}
	
	\begin{frame}[t]{Q391 -- Answer \& Explanation}
		\small
		\textbf{Correct Answer: B} \\[4pt]
		\textbf{Why:} Standard 1312 requires the external assessor to be independent and objective. Consulting services to the function being assessed create a self-review threat that most limits the assessment's credibility. \\[4pt]
		\textbf{A wrong:} Qualifications are presumed adequate. \\[2pt]
		\textbf{C wrong:} Cost is not the key concern. \\[2pt]
		\textbf{D wrong:} Duration is not the key concern.
	\end{frame}
	
	% ============================================================
	% Q392
	% ============================================================
	\begin{frame}[t]{Q392}
		\small
		\textbf{Question:} The CAE is evaluating whether the QAIP includes all required elements. The QAIP includes internal assessments and external assessments but does not include a process for monitoring whether recommendations from assessments are implemented. Which is the MOST accurate statement?
		\begin{enumerate}[label=\Alph*.,itemsep=1pt,topsep=2pt]
			\item The QAIP is complete because internal and external assessments exist.
			\item The QAIP is incomplete; a QAIP must include monitoring of whether recommendations are implemented and reporting results to the board (Standards 1300/1320).
			\item Monitoring recommendations is optional.
			\item The external assessment covers monitoring.
		\end{enumerate}
	\end{frame}
	
	\begin{frame}[t]{Q392 -- Answer \& Explanation}
		\small
		\textbf{Correct Answer: B} \\[4pt]
		\textbf{Why:} Standard 1300 requires a QAIP covering all aspects of the activity, including follow-up on assessment recommendations and reporting to the board (Standard 1320). \\[4pt]
		\textbf{A wrong:} Assessments without follow-up are incomplete. \\[2pt]
		\textbf{C wrong:} Monitoring recommendations is required. \\[2pt]
		\textbf{D wrong:} External assessment does not cover ongoing monitoring.
	\end{frame}
	
	% ============================================================
	% Q393
	% ============================================================
	\begin{frame}[t]{Q393}
		\small
		\textbf{Question:} The CAE is asked to evaluate whether the internal audit activity conforms with the Standards. The CAE finds that the activity's charter is approved by the CEO, not the board, and that the audit committee has never seen the charter. Which is the MOST accurate statement?
		\begin{enumerate}[label=\Alph*.,itemsep=1pt,topsep=2pt]
			\item The activity conforms because the CEO is senior management.
			\item The activity does not conform; Standard 1000 requires the charter to be approved by the board (or its audit committee), and the CAE must disclose the nonconformance and seek remediation.
			\item The activity conforms because the audit committee approves the plan.
			\item The activity conforms because the charter exists.
		\end{enumerate}
	\end{frame}
	
	\begin{frame}[t]{Q393 -- Answer \& Explanation}
		\small
		\textbf{Correct Answer: B} \\[4pt]
		\textbf{Why:} Standard 1000 requires board approval of the charter. CEO approval with no audit committee involvement is a nonconformance requiring disclosure and remediation (Standard 1320). \\[4pt]
		\textbf{A wrong:} Management approval does not satisfy Standard 1000. \\[2pt]
		\textbf{C wrong:} Plan approval does not substitute for charter approval. \\[2pt]
		\textbf{D wrong:} Existence is not approval.
	\end{frame}
	
	% ============================================================
	% Q394
	% ============================================================
	\begin{frame}[t]{Q394}
		\small
		\textbf{Question:} The CAE is evaluating the QAIP's periodic internal assessments. The activity performed a self-assessment using a recognized maturity model but did not validate the results independently. Which is the MOST accurate statement?
		\begin{enumerate}[label=\Alph*.,itemsep=1pt,topsep=2pt]
			\item The assessment is sufficient because a recognized model was used.
			\item The assessment should be validated or supplemented by independent review to support the conclusion; self-assessment alone may lack objectivity.
			\item Periodic internal assessments are not required.
			\item The external assessment alone is sufficient.
		\end{enumerate}
	\end{frame}
	
	\begin{frame}[t]{Q394 -- Answer \& Explanation}
		\small
		\textbf{Correct Answer: B} \\[4pt]
		\textbf{Why:} Standard 1311 requires periodic internal assessments; using a recognized model is good, but objectivity is enhanced by independent validation. Self-assessment alone may lack sufficient objectivity for the conclusion. \\[4pt]
		\textbf{A wrong:} A model alone does not ensure objectivity. \\[2pt]
		\textbf{C wrong:} Periodic internal assessments are required. \\[2pt]
		\textbf{D wrong:} External assessment does not replace periodic internal assessments.
	\end{frame}
	
	% ============================================================
	% Q395
	% ============================================================
	\begin{frame}[t]{Q395}
		\small
		\textbf{Question:} The CAE is evaluating the internal audit activity's conformance. Findings: (1) QAIP exists; (2) external assessment within five years; (3) one engagement's workpapers were destroyed 2 months after the report was issued; (4) the plan is risk-based. Which is the MOST accurate statement?
		\begin{enumerate}[label=\Alph*.,itemsep=1pt,topsep=2pt]
			\item The activity conforms because the QAIP and external assessment are in place.
			\item The activity does not fully conform; Standards 2330/2330.A1 require retention of engagement records per policy and legal requirements, and destroying workpapers 2 months after issuance likely violates retention requirements.
			\item Destruction of workpapers is not a conformance issue.
			\item The risk-based plan cures the retention issue.
		\end{enumerate}
	\end{frame}
	
	\begin{frame}[t]{Q395 -- Answer \& Explanation}
		\small
		\textbf{Correct Answer: B} \\[4pt]
		\textbf{Why:} Standards 2330/2330.A1 require retention of engagement records per organizational policy and legal requirements. Destroying workpapers 2 months after issuance likely violates retention and is a nonconformance. \\[4pt]
		\textbf{A wrong:} QAIP/external assessment do not cure the retention issue. \\[2pt]
		\textbf{C wrong:} Retention is a conformance requirement. \\[2pt]
		\textbf{D wrong:} The plan does not cure retention.
	\end{frame}
	
	% ============================================================
	% Q396
	% ============================================================
	\begin{frame}[t]{Q396}
		\small
		\textbf{Question:} The CAE is asked by the audit committee whether the internal audit activity ``conforms'' with the Standards. The activity has one significant nonconformance: the CAE approved an engagement that audited a process the CAE designed two years ago. The nonconformance was disclosed to the board and a remediation plan was implemented. Which is the MOST accurate statement?
		\begin{enumerate}[label=\Alph*.,itemsep=1pt,topsep=2pt]
			\item The activity conforms because the nonconformance was disclosed and remediated.
			\item The activity does not currently conform with Standard 1130.A1 regarding that engagement, but the disclosure and remediation satisfy Standard 1320; the CAE should describe the nonconformance, its impact, and the remediation plan.
			\item The activity conforms because two years have passed.
			\item The activity conforms because the board accepted the remediation plan.
		\end{enumerate}
	\end{frame}
	
	\begin{frame}[t]{Q396 -- Answer \& Explanation}
		\small
		\textbf{Correct Answer: B} \\[4pt]
		\textbf{Why:} A self-review engagement (auditing a process the CAE designed) is a nonconformance with Standard 1130.A1. Disclosure and remediation satisfy Standard 1320, but the CAE must accurately describe the nonconformance, its impact, and the plan. \\[4pt]
		\textbf{A wrong:} Remediation does not erase the nonconformance. \\[2pt]
		\textbf{C wrong:} Time does not cure self-review. \\[2pt]
		\textbf{D wrong:} Acceptance of a plan does not erase the nonconformance.
	\end{frame}
	
	% ============================================================
	% Q397
	% ============================================================
	\begin{frame}[t]{Q397}
		\small
		\textbf{Question:} The internal audit activity has a QAIP, but the QAIP is not documented. Which is the MOST accurate statement?
		\begin{enumerate}[label=\Alph*.,itemsep=1pt,topsep=2pt]
			\item The QAIP conforms because the activities are performed.
			\item The QAIP does not fully conform; documentation of the QAIP and its results is required (Standards 1300/1320).
			\item Documentation of the QAIP is optional.
			\item The external assessment documents the QAIP.
		\end{enumerate}
	\end{frame}
	
	\begin{frame}[t]{Q397 -- Answer \& Explanation}
		\small
		\textbf{Correct Answer: B} \\[4pt]
		\textbf{Why:} Standard 1300/1320 require a documented QAIP with results reported to the board. Undocumented activities cannot be evidenced. \\[4pt]
		\textbf{A wrong:} Documentation is required. \\[2pt]
		\textbf{C wrong:} Documentation is not optional. \\[2pt]
		\textbf{D wrong:} External assessment does not document the QAIP.
	\end{frame}
	
	% ============================================================
	% Q398
	% ============================================================
	\begin{frame}[t]{Q398}
		\small
		\textbf{Question:} The CAE is evaluating whether the internal audit activity can conclude ``conforms'' when the external assessment concluded ``generally conforms'' with one ``partially conforms'' finding related to engagement supervision. Which is the MOST accurate statement?
		\begin{enumerate}[label=\Alph*.,itemsep=1pt,topsep=2pt]
			\item The activity can conclude ``conforms'' because the overall conclusion was ``generally conforms.''
			\item The ``partially conforms'' finding on supervision indicates a nonconformance; the CAE must evaluate its impact, disclose it to the board if significant, and remediate (Standards 1320/2340).
			\item ``Partially conforms'' is not a nonconformance.
			\item The external assessor's conclusion is binding on the CAE.
		\end{enumerate}
	\end{frame}
	
	\begin{frame}[t]{Q398 -- Answer \& Explanation}
		\small
		\textbf{Correct Answer: B} \\[4pt]
		\textbf{Why:} ``Partially conforms'' indicates a nonconformance (Standard 2340 supervision). The CAE must evaluate impact, disclose if significant, and remediate (Standard 1320). \\[4pt]
		\textbf{A wrong:} ``Generally conforms'' with a partial nonconformance is not full conformance. \\[2pt]
		\textbf{C wrong:} It is a nonconformance. \\[2pt]
		\textbf{D wrong:} The CAE evaluates the conclusion; it is not blindly binding.
	\end{enumerate}
	
	% ============================================================
	% Q399
	% ============================================================
	\begin{frame}[t]{Q399}
		\small
		\textbf{Question:} The CAE is evaluating the QAIP's coverage of the internal audit activity. The QAIP covers engagement performance but not the activity's governance, administration, or resource management. Which is the MOST accurate statement?
		\begin{enumerate}[label=\Alph*.,itemsep=1pt,topsep=2pt]
			\item The QAIP is adequate because engagement performance is the core of internal audit.
			\item The QAIP is incomplete; it must cover all aspects of the internal audit activity, including governance, administration, and resource management (Standard 1300).
			\item Governance and administration are outside the QAIP.
			\item Resource management is covered by the budget.
		\end{enumerate}
	\end{frame}
	
	\begin{frame}[t]{Q399 -- Answer \& Explanation}
		\small
		\textbf{Correct Answer: B} \\[4pt]
		\textbf{Why:} Standard 1300 requires the QAIP to cover all aspects of the internal audit activity, not just engagements. \\[4pt]
		\textbf{A wrong:} Engagement performance alone is insufficient. \\[2pt]
		\textbf{C wrong:} Governance and administration are within the QAIP. \\[2pt]
		\textbf{D wrong:} The budget does not substitute for QAIP coverage.
	\end{frame}
	
	% ============================================================
	% Q400
	% ============================================================
	\begin{frame}[t]{Q400}
		\small
		\textbf{Question:} The CAE is asked by the audit committee whether the external quality assessment can be performed by the organization's external auditor. Which is the MOST accurate statement?
		\begin{enumerate}[label=\Alph*.,itemsep=1pt,topsep=2pt]
			\item The external auditor cannot perform the assessment because it would impair the auditor's independence.
			\item The external auditor may perform the assessment if qualified and independent of the internal audit activity; the CAE must evaluate the auditor's objectivity relative to internal audit.
			\item The external auditor is the only permitted assessor.
			\item The external auditor's assessment is not valid.
		\end{enumerate}
	\end{frame}
	
	\begin{frame}[t]{Q400 -- Answer \& Explanation}
		\small
		\textbf{Correct Answer: B} \\[4pt]
		\textbf{Why:} Standard 1312 permits a qualified, independent external assessor. The external auditor may qualify if independent of the internal audit activity; the CAE must evaluate objectivity relative to internal audit. \\[4pt]
		\textbf{A wrong:} External auditor independence from the organization does not automatically preclude the assessment. \\[2pt]
		\textbf{C wrong:} Other qualified assessors are permitted. \\[2pt]
		\textbf{D wrong:} A valid assessment by the external auditor is possible.
	\end{frame}
	
	% ============================================================
	% Q401
	% ============================================================
	\begin{frame}[t]{Q401}
		\small
		\textbf{Question:} The CAE is evaluating the internal audit activity's conformance when the QAIP was implemented only 6 months ago and no external assessment has been performed. Which is the MOST accurate statement?
		\begin{enumerate}[label=\Alph*.,itemsep=1pt,topsep=2pt]
			\item The activity conforms because the QAIP was recently implemented.
			\item The activity does not fully conform; the external assessment requirement applies (Standard 1312), and the CAE should plan for it and disclose the current status to the board.
			\item External assessment is optional for newly formed activities.
			\item The QAIP implementation satisfies the external assessment requirement.
		\end{enumerate}
	\end{frame}
	
	\begin{frame}[t]{Q401 -- Answer \& Explanation}
		\small
		\textbf{Correct Answer: B} \\[4pt]
		\textbf{Why:} Standard 1312 requires external assessment at least every five years; the requirement is not waived for newly formed activities. The CAE should plan for it and disclose current status to the board. \\[4pt]
		\textbf{A wrong:} Recency does not waive the requirement. \\[2pt]
		\textbf{C wrong:} External assessment is not optional. \\[2pt]
		\textbf{D wrong:} QAIP implementation does not satisfy the external assessment.
	\end{frame}
	
	% ============================================================
	% Q402
	% ============================================================
	\begin{frame}[t]{Q402}
		\small
		\textbf{Question:} The CAE is evaluating the QAIP's reporting to the board. The QAIP results are reported to the CEO annually but not to the audit committee. Which is the MOST accurate statement?
		\begin{enumerate}[label=\Alph*.,itemsep=1pt,topsep=2pt]
			\item Reporting to the CEO is sufficient.
			\item Reporting to the CEO is insufficient; Standards 1320/2430 require reporting QAIP results and conformance to the board/audit committee.
			\item QAIP results are internal only.
			\item The external assessor reports to the board on the CAE's behalf.
		\end{enumerate}
	\end{frame}
	
	\begin{frame}[t]{Q402 -- Answer \& Explanation}
		\small
		\textbf{Correct Answer: B} \\[4pt]
		\textbf{Why:} Standards 1320/2430 require the CAE to report QAIP results and conformance with the Standards to the board/audit committee. CEO-only reporting is insufficient. \\[4pt]
		\textbf{A wrong:} CEO reporting does not satisfy board reporting. \\[2pt]
		\textbf{C wrong:} QAIP results must be reported to the board. \\[2pt]
		\textbf{D wrong:} The CAE reports; the external assessor does not act on the CAE's behalf.
	\end{frame}
	
	% ============================================================
	% Q403
	% ============================================================
	\begin{frame}[t]{Q403}
		\small
		\textbf{Question:} The CAE is evaluating the internal audit activity's conformance when an engagement was performed by a co-sourced provider whose work was not reviewed by the CAE before the report was issued. Which is the MOST accurate statement?
		\begin{enumerate}[label=\Alph*.,itemsep=1pt,topsep=2pt]
			\item The activity conforms because the provider is qualified.
			\item The activity does not fully conform; the CAE must ensure adequate supervision and review of co-sourced work before issuance (Standards 2340/2070), and the lack of review is a nonconformance.
			\item Co-sourced work does not require review.
			\item The provider's own QAIP covers the CAE's responsibility.
		\end{enumerate}
	\end{frame}
	
	\begin{frame}[t]{Q403 -- Answer \& Explanation}
		\small
		\textbf{Correct Answer: B} \\[4pt]
		\textbf{Why:} Standards 2340 and 2070 require adequate supervision and review of engagements, including co-sourced work, before reports are issued. The CAE retains responsibility. \\[4pt]
		\textbf{A wrong:} Provider qualifications do not eliminate the CAE's review responsibility. \\[2pt]
		\textbf{C wrong:} Co-sourced work requires review. \\[2pt]
		\textbf{D wrong:} The provider's QAIP does not cover the CAE's responsibility.
	\end{frame}
	
	% ============================================================
	% Q404
	% ============================================================
	\begin{frame}[t]{Q404}
		\small
		\textbf{Question:} The CAE is evaluating the internal audit activity's conformance when the QAIP identified a significant nonconformance in engagement planning, but the CAE did not report it to the board. Which is the MOST accurate statement?
		\begin{enumerate}[label=\Alph*.,itemsep=1pt,topsep=2pt]
			\item The nonconformance is a QAIP matter only and does not require board reporting.
			\item The CAE violated Standard 1320 by not reporting the significant nonconformance to senior management and the board.
			\item The nonconformance should be reported only if it affects the annual opinion.
			\item The external assessor is responsible for reporting it.
		\end{enumerate}
	\end{frame}
	
	\begin{frame}[t]{Q404 -- Answer \& Explanation}
		\small
		\textbf{Correct Answer: B} \\[4pt]
		\textbf{Why:} Standard 1320 requires the CAE to report significant nonconformance to senior management and the board. Non-disclosure is itself a violation. \\[4pt]
		\textbf{A wrong:} Significant nonconformance requires board reporting. \\[2pt]
		\textbf{C wrong:} Reporting is not limited to annual opinion impact. \\[2pt]
		\textbf{D wrong:} The CAE is responsible, not the external assessor.
	\end{frame}
	
	% ============================================================
	% Q405
	% ============================================================
	\begin{frame}[t]{Q405}
		\small
		\textbf{Question:} The CAE is evaluating whether the internal audit activity can conclude ``conforms'' when the external assessment is 4 years and 11 months old, and the QAIP shows only minor documentation gaps in engagement workpapers. Which is the MOST accurate statement?
		\begin{enumerate}[label=\Alph*.,itemsep=1pt,topsep=2pt]
			\item The activity cannot conclude ``conforms'' because the external assessment is about to expire.
			\item The external assessment is within the five-year requirement; the minor documentation gaps should be evaluated for significance and disclosed if significant, but they do not automatically preclude ``conforms.''
			\item The minor gaps automatically preclude ``conforms.''
			\item The external assessment must be repeated immediately.
		\end{enumerate}
	\end{frame}
	
	\begin{frame}[t]{Q405 -- Answer \& Explanation}
		\small
		\textbf{Correct Answer: B} \\[4pt]
		\textbf{Why:} 4 years 11 months is within the five-year requirement (Standard 1312). Minor documentation gaps are nonconformances only if significant; the CAE evaluates significance and discloses if required (Standard 1320). \\[4pt]
		\textbf{A wrong:} The assessment is still within the window. \\[2pt]
		\textbf{C wrong:} Minor gaps do not automatically preclude conformance. \\[2pt]
		\textbf{D wrong:} Immediate repetition is not required.
	\end{frame}
	
	% ============================================================
	% Q406
	% ============================================================
	\begin{frame}[t]{Q406}
		\small
		\textbf{Question:} The CAE is evaluating the internal audit activity's conformance when the QAIP is designed but not yet operating (no assessments performed). Which is the MOST accurate statement?
		\begin{enumerate}[label=\Alph*.,itemsep=1pt,topsep=2pt]
			\item The activity conforms because the QAIP is designed.
			\item The activity does not fully conform; the QAIP must be operating, not merely designed (Standard 1300).
			\item Design is sufficient for conformance.
			\item The external assessment will satisfy the operating requirement.
		\end{enumerate}
	\end{frame}
	
	\begin{frame}[t]{Q406 -- Answer \& Explanation}
		\small
		\textbf{Correct Answer: B} \\[4pt]
		\textbf{Why:} Standard 1300 requires the QAIP to be operating (assessments performed), not merely designed. \\[4pt]
		\textbf{A wrong:} Design without operation is insufficient. \\[2pt]
		\textbf{C wrong:} Design is not sufficient. \\[2pt]
		\textbf{D wrong:} External assessment does not cover ongoing QAIP operation.
	\end{enumerate}
	
	% ============================================================
	% Q407
	% ============================================================
	\begin{frame}[t]{Q407}
		\small
		\textbf{Question:} The CAE is evaluating the internal audit activity's conformance when the QAIP includes ongoing and periodic internal assessments and external assessment, but the CAE has not reported the QAIP results to the board in the last two years. Which is the MOST accurate statement?
		\begin{enumerate}[label=\Alph*.,itemsep=1pt,topsep=2pt]
			\item The activity conforms because the QAIP elements exist.
			\item The activity does not fully conform; Standards 1320/2430 require reporting QAIP results and conformance to the board.
			\item Reporting QAIP results is optional.
			\item The external assessor's report substitutes for CAE reporting.
		\end{enumerate}
	\end{frame}
	
	\begin{frame}[t]{Q407 -- Answer \& Explanation}
		\small
		\textbf{Correct Answer: B} \\[4pt]
		\textbf{Why:} Standards 1320/2430 require the CAE to report QAIP results and conformance to the board. Two years without reporting is a nonconformance. \\[4pt]
		\textbf{A wrong:} Existence of elements without reporting is insufficient. \\[2pt]
		\textbf{C wrong:} Reporting is required. \\[2pt]
		\textbf{D wrong:} The external assessor does not substitute for CAE reporting.
	\end{enumerate}
	
	% ============================================================
	% Q408
	% ============================================================
	\begin{frame}[t]{Q408}
		\small
		\textbf{Question:} The CAE is evaluating whether the internal audit activity conforms when an engagement was performed without a documented work program. Which is the MOST accurate statement?
		\begin{enumerate}[label=\Alph*.,itemsep=1pt,topsep=2pt]
			\item The activity conforms because the engagement was completed.
			\item The activity does not conform; Standard 2240 requires a documented work program for each engagement.
			\item Work programs are optional if the auditor is experienced.
			\item Work programs are required only for assurance engagements.
		\end{enumerate}
	\end{frame}
	
	\begin{frame}[t]{Q408 -- Answer \& Explanation}
		\small
		\textbf{Correct Answer: B} \\[4pt]
		\textbf{Why:} Standard 2240 requires a documented work program for each engagement (assurance and consulting). \\[4pt]
		\textbf{A wrong:} Completion without documentation is a nonconformance. \\[2pt]
		\textbf{C wrong:} Experience does not waive the requirement. \\[2pt]
		\textbf{D wrong:} Work programs are required for all engagements.
	\end{enumerate}
	
	% ============================================================
	% Q409
	% ============================================================
	\begin{frame}[t]{Q409}
		\small
		\textbf{Question:} The CAE is evaluating whether the internal audit activity conforms when engagement conclusions are not supported by sufficient evidence in the workpapers. Which is the MOST accurate statement?
		\begin{enumerate}[label=\Alph*.,itemsep=1pt,topsep=2pt]
			\item The activity conforms because conclusions were reached.
			\item The activity does not conform; Standards 2300/2320/2330 require sufficient, reliable, relevant, useful evidence and documentation supporting conclusions.
			\item Evidence sufficiency is a judgment that cannot be assessed.
			\item The report's conclusions do not need workpaper support.
		\end{enumerate}
	\end{frame}
	
	\begin{frame}[t]{Q409 -- Answer \& Explanation}
		\small
		\textbf{Correct Answer: B} \\[4pt]
		\textbf{Why:} Standards 2300/2320/2330 require sufficient evidence and documentation supporting conclusions. Unsupported conclusions are a nonconformance. \\[4pt]
		\textbf{A wrong:} Conclusions without evidence are improper. \\[2pt]
		\textbf{C wrong:} Sufficiency is assessable against the Standards. \\[2pt]
		\textbf{D wrong:} Conclusions require workpaper support.
	\end{enumerate}
	
	% ============================================================
	% Q410
	% ============================================================
	\begin{frame}[t]{Q410}
		\small
		\textbf{Question:} The CAE is evaluating whether the internal audit activity conforms when the QAIP includes a maturity model self-assessment, but the results were not validated and the CAE concluded ``conforms'' without further evidence. Which is the MOST accurate statement?
		\begin{enumerate}[label=\Alph*.,itemsep=1pt,topsep=2pt]
			\item The conclusion ``conforms'' is adequately supported by the self-assessment.
			\item The conclusion is not adequately supported; self-assessment alone may lack objectivity and the CAE should obtain additional evidence (e.g., independent review) before concluding ``conforms.''
			\item Maturity models are not acceptable for QAIP.
			\item The board's acceptance of the conclusion cures the evidence gap.
		\end{enumerate}
	\end{frame}
	
	\begin{frame}[t]{Q410 -- Answer \& Explanation}
		\small
		\textbf{Correct Answer: B} \\[4pt]
		\textbf{Why:} A self-assessment alone may lack objectivity; the CAE needs sufficient evidence (e.g., independent review or external assessment) to support a ``conforms'' conclusion. \\[4pt]
		\textbf{A wrong:} Self-assessment alone is insufficient. \\[2pt]
		\textbf{C wrong:} Maturity models are acceptable, but not alone. \\[2pt]
		\textbf{D wrong:} Board acceptance does not cure an evidence gap.
	\end{enumerate}
	
	% ============================================================
	% Q411
	% ============================================================
	\begin{frame}[t]{Q411}
		\small
		\textbf{Question:} The CAE is evaluating whether the internal audit activity conforms when the external assessment was performed by a team that included a former internal audit staff member who left the organization 8 months ago. Which is the MOST accurate statement?
		\begin{enumerate}[label=\Alph*.,itemsep=1pt,topsep=2pt]
			\item The assessment conforms because the team is otherwise qualified.
			\item The assessment's independence is impaired; a former staff member who left 8 months ago creates an independence threat, and the CAE should evaluate and likely obtain a different assessor.
			\item Former staff members are always prohibited from external assessments.
			\item The organization's independence from the assessor cures the issue.
		\end{enumerate}
	\end{frame}
	
	\begin{frame}[t]{Q411 -- Answer \& Explanation}
		\small
		\textbf{Correct Answer: B} \\[4pt]
		\textbf{Why:} Standard 1312 requires the external assessor to be independent and objective. A former staff member who left 8 months ago creates an independence threat (self-review/familiarity); the CAE should evaluate and likely replace. \\[4pt]
		\textbf{A wrong:} Team qualifications do not cure the independence threat. \\[2pt]
		\textbf{C wrong:} Not always prohibited, but the threat must be evaluated. \\[2pt]
		\textbf{D wrong:} Organizational independence does not cure individual threats.
	\end{enumerate}
	
	% ============================================================
	% Q412
	% ============================================================
	\begin{frame}[t]{Q412}
		\small
		\textbf{Question:} The CAE is evaluating whether the internal audit activity conforms when the QAIP includes ongoing assessments but the CAE's annual report to the board does not mention conformance with the Standards. Which is the MOST accurate statement?
		\begin{enumerate}[label=\Alph*.,itemsep=1pt,topsep=2pt]
			\item The activity conforms because ongoing assessments exist.
			\item The activity does not fully conform; Standard 2430 requires the CAE to report conformance with the Standards and Code to the board at least annually.
			\item Conformance reporting is optional.
			\item The external assessor's report substitutes for the annual conformance report.
		\end{enumerate}
	\end{frame}
	
	\begin{frame}[t]{Q412 -- Answer \& Explanation}
		\small
		\textbf{Correct Answer: B} \\[4pt]
		\textbf{Why:} Standard 2430 requires the CAE to report conformance with the Standards and Code to the board at least annually. Omitting it is a nonconformance. \\[4pt]
		\textbf{A wrong:} Ongoing assessments do not substitute for conformance reporting. \\[2pt]
		\textbf{C wrong:} Conformance reporting is required. \\[2pt]
		\textbf{D wrong:} External assessor reporting does not substitute.
	\end{enumerate}
	
	% ============================================================
	% Q413
	% ============================================================
	\begin{frame}[t]{Q413}
		\small
		\textbf{Question:} The CAE is asked to conclude on conformance. The QAIP identified three nonconformances: (1) one engagement audited a process the partner designed; (2) workpapers for one engagement were destroyed early; (3) the external assessment is 3 years old. Which statement is MOST accurate regarding the conformance conclusion?
		\begin{enumerate}[label=\Alph*.,itemsep=1pt,topsep=2pt]
			\item The activity conforms because the external assessment is within 5 years and only two nonconformances exist.
			\item The activity does not conform; the self-review engagement and premature destruction are nonconformances that must be evaluated and reported, and the CAE should describe impact and remediation (Standard 1320).
			\item Only the self-review engagement matters.
			\item The 3-year-old external assessment is a nonconformance.
		\end{enumerate}
	\end{frame}
	
	\begin{frame}[t]{Q413 -- Answer \& Explanation}
		\small
		\textbf{Correct Answer: B} \\[4pt]
		\textbf{Why:} The self-review engagement (Standard 1130.A1) and premature destruction of workpapers (Standard 2330.A1) are nonconformances requiring evaluation, disclosure to the board, and remediation (Standard 1320). The 3-year-old external assessment is within the five-year window. \\[4pt]
		\textbf{A wrong:} Multiple nonconformances preclude ``conforms.'' \\[2pt]
		\textbf{C wrong:} Both the self-review and retention issues matter. \\[2pt]
		\textbf{D wrong:} 3 years is within the five-year requirement.
	\end{frame}
	
	\begin{frame}{End}
		\centering
		\Large End of Questions 355--413\\[8pt]
		\normalsize Domains: CAE--Board Relations \& Independence Threats; QAIP, Conformance Judgment \& External Assessment.
	\end{frame}
	
\end{document}